\begin{afterword}
\summaryhead

The post asks what kind of humility good science needs. It sorts cases. The creationist who urges humility about evolution practices selective underconfidence. The engineer who builds fail-safes while sure the machine will not fail, and the student who double-checks a test, show good humility: they act against their own possible errors. The student who will not check because nothing is certain is less wise, and the student who says ``one so lowly as myself'' cannot improve is showing social modesty, which concerns status, not scientific humility, which you would practice alone in space.

Vague ideas of humility, the post argues, can justify anything, and most appeals to humility are excuses to shrug, a form of motivated skepticism. Humility is also easy to profess while changing nothing, as Daniel Dennett says of religious ``belief in belief.'' The test is whether the actions a humble line of thought recommends make you stronger or weaker. The post ends with a definition: ``To be humble is to take specific actions in anticipation of your own errors.''

\responsehead

The post draws a useful line between humility that changes what you do and humility that is only professed, and its test, to judge a humble thought by the actions it recommends, is clear. The problem is the word it attaches to that line.

In ordinary use, humility is an attitude to oneself: Merriam-Webster gives ``freedom from pride or arrogance,'' and for humble ``not proud or haughty'' and ``ranking low in a hierarchy or scale.'' The post knows this. It says that people asked to ``be more humble'' will think of social modesty, and it calls its own sense ``scientific humility.'' But it then denies the everyday sense the name (``This is social modesty, not humility'') and gives the word to double-checking and fail-safes, which are closer to what the dictionary calls prudence. Pride drops out of it. The post's model engineer is ``damn sure.'' On the post's definition a proud person can be humble. The same definition, quoted here from ``Twelve Virtues of Rationality,'' comes back later in the book, and our notes there make this point again.

The account of where social modesty comes from is asserted. That it is ``ancestrally relevant'' and that scientific humility is ``a more recent and rarefied invention'' are offered without evidence. They are the post's whole account of why the everyday sense belongs to status and not to science.

The post's support is mostly its own examples and observation, as a short essay's usually is, and it hedges them as such (``The vast majority of appeals that I witness''). Its one cited study is described accurately. The student who raises the problem of induction first gets a motive for an answer, ``go home and play video games,'' and gets a reasoned reply only four paragraphs later, in the bridge test.

Twice the post introduces a sentence with ``Therefore it is written.'' Both come from the author's own ``Twelve Virtues of Rationality,'' and the post does not say so.

In short: a useful distinction between humility practiced and humility professed, drawn under a definition of humility that leaves out freedom from pride.
\end{afterword}
